\documentclass{article}
\usepackage[utf8]{inputenc}
\usepackage{amsmath, hyperref}

\pagenumbering{gobble}

\title{Year 8 Maths Worksheets}
\author{C C Barnes}
\date{September 2022}

\begin{document}

\maketitle

\begin{abstract}
    A (very small) collection of Year 8 worksheets seeing as I only taught this year once. As with most of my worksheets, I kept everything for a year group in a single .tex file and only printed the relevant pages as needed. Any title pages for topics are purely for organization and serve no further purpose.
\end{abstract}

\tableofcontents

\newpage

\section{Standard Form}

\subsection{Solutions}

\begin{enumerate}
    \item[a)] \(6 \times 10^{3+n}\)
    \item[b)] \(1 \times 10^{n-2}\)
    \item[c)] \(6 \times 10^{2n}\)
    \item[d)] \(1 \times 10^{2n+1}\)
    \item[e)] \(4 \times 10^{2n}\)
    \item[f)] \(5 \times 10^{2n-1}\)
    \item[g)] \(9 \times 10^{2n}\)
    \item[h)] \(3 \times 10^n\)
\end{enumerate}

\begin{enumerate}
    \item [4)] \begin{enumerate}
        \item[a)] \(\frac{1}{4}\leq x < 2\frac{1}{2}\)
        \item[b)] \(10\leq x < 100\)
        \item[c)] \(0\leq x<9\)
        \item[d)] \(0\leq x < 4\frac{1}{2}\)
        \item[e)] \(\frac{3}{4}\leq x < 3\)
        \item[f)] \(1\leq x < \sqrt{10}\)
        \item[g)] \(1 \leq x < 100\)
    \end{enumerate}
\end{enumerate}

\newpage
\section{Pythagoras' Theorem}

\newpage

\noindent \textbf{Coordinate Geometry using Pythagoras' Theorem}
\medskip
\hrule
\medskip

\noindent For each of these questions, leave your answers as a surd.

\begin{enumerate}
    \item What is the distance between the coordinates \begin{enumerate}
        \item (3,5) and (7,4)
        \item (-1,2) and (3,6)
        \item (-3, -5) and (-1, 7)
        \item (3, -4) and (6, -2)
    \end{enumerate}
    \item Find a coordinate with integer values that is
    \begin{enumerate}
        \item a distance of \(\sqrt{5}\) from (1,2)
        \item a distance of \(\sqrt{13}\) from (-4,5)
        \item a distance of \(\sqrt{53}\) from (-3, -7)
        \item a distance of \(\sqrt{40}\) from (3,-5)
    \end{enumerate}
    \item A right angled isosceles triangle has area equal to 18 cm\(^2\). What are the lengths of all 3 sides?
\end{enumerate}

\noindent \textbf{Extension}

\medskip

These are more challenging questions to be attempted only by the bravest of mathematicians. You are welcome to discuss in pairs or triplets (instead of pairs or triplets we ought to say as mathematicians \(n\)-tuples where \(n \in \{2,3\}\)).

\begin{enumerate}
    \item Consider again question 2 above. Suppose we remove the restriction that our coordinate should have integer values. There are now infinitely many possible answers but they form a 2D-shape. Which shape is this?
    \item Suppose we now go to three-dimensional coordinates. Name a coordinate with integer values that is a distance of \(\sqrt{3}\) from (0,0,0).
    \item Again, drop the requirement of having integer solutions only. Thus there are infinitely many points that lie a distance of \(\sqrt{3}\) from the point \((0,0,0)\). What shape is this?
\end{enumerate}

\newpage

\section{Straight Line Graphs}
\newpage 
\noindent \textbf{Finding Gradients and \(y\)-intercepts from Equations}
\medskip
\hrule
\medskip
\noindent \textbf{Extension} 

\medskip

Find the exact value of the gradient and \(y\)-intercept of the following (and yes, they are all in fact lines despite their strange appearance).

\begin{enumerate}
    \item \(y-3 = \frac{2}{7}(3x-4)\)
    \item \(y^2-7x+2y = 64 - 3x +y^2 \)
    \item \(2z + 3x^2z - 4\pi y = z(2+3x^2) +3y - 2x -\sqrt{2}\)
\end{enumerate}

% Why not generalise the coordinate axes. So, for instance, have a line in the p-q plane or a-b plane etc.

% Why not add in arbitrary constants. 

\newpage

\noindent \textbf{Answers}

\medskip

\hrule

\medskip

\begin{enumerate}
    \item \begin{enumerate}
        \item grad = \(\frac{6}{7}\)
        \item \(y\)-intercept \(= \frac{13}{7}\)
    \end{enumerate}
    \item \begin{enumerate}
        \item grad = \(2\)
        \item \(y\)-intercept = 32
    \end{enumerate}
    \item \begin{enumerate}
        \item grad = \(\frac{2}{3+4\pi}\)
        \item \(y\)-intercept = \(\frac{\sqrt{2}}{3+4\pi}\)
    \end{enumerate}
\end{enumerate}

\newpage

\noindent \textbf{Interpreting Gradients and Intercepts}

\medskip

\hrule

\medskip

\begin{enumerate}
    \item A study was done to see the relationship between the time it takes, \(y\), measured in years, to complete a university degree and the student loan debt incurred, \(x\), measured in £. It was found that
    \[y=25142+14329x\]
    Interpret the gradient of this equation in the context of the study.
    \item Suppose that a research group tested the cholesterol level of a sample of 40 year old women and then waited many years to see the relationship between a woman's HDL cholesterol level in mg/dl,  \(x\), and her age at death,  \(y\). The equation of the regression line was found to be:
    \[y=103 - 0.3x\]
    Interpret the slope of the regression line in the context of the study.
    \item A ball is thrown into the air and its speed is given by the linear graph
    \[v=128-6t\]
    \begin{enumerate}
        \item What is the gradient in this context?
        \item What is the intercept?
        \item What are the units of the gradient?
        \item What are the units of the \(y\)-intercept?
    \end{enumerate}
    \item The injuries, \(y\), of track and field athletes is found to correlate with the number of hours spent training, \(x\), where \(x\) is measured in thousands of hours. It is found that
    \[
    y = 2+2.64x.
    \]
    What is the meaning of the gradient?
    \item A recent study found a relationship between salt intake in the diet, \(x\), measured in g, and systolic blood pressure, \(y\), measured in millimetres of mercury (mmHg) to be
    \[y = 10x + 60\]
    Interpret the gradient in this relationship. What are the units?
\end{enumerate}


\newpage

\noindent \textbf{Answers}

\medskip

\hrule

\medskip

\begin{enumerate}
    \item The cost per additional year of studying in £.
    \item The loss in lifespan per mg/dl increase in cholesterol.
    \item \begin{enumerate}
        \item The change in velocity per unit time or acceleration.
        \item The initial velocity when the ball is thrown.
        \item ms\(^{-2}\)
        \item ms\(^{-1}\)
    \end{enumerate}
    \item The gradient is the additional injuries for every additional thousand hours of practice.
    \item The increase in blood pressure per daily 1g increase in salt intake. The units are mmHg/g.
    
\end{enumerate}

\end{document}
